% tex/impl/commands/graphs.tex
% Wiederverwendbare Notation fuer gerichtete und ungerichtete Graphen.

\ProvidesFile{graphs.tex}[2026/09/02 Graph-theoretic notation]

% Strukturpraedikate. Wie bei \PartOrd erhaelt jedes Makro die gesamte
% Parameterliste als ein Argument; geschrieben wird also etwa
% \GraphStruct{V,E}.
\DeclareRobustCommand{\GraphStruct}[1]{%
  \ensuremath{\mathsf{Graph}\!\left(#1\right)}%
}
\DeclareRobustCommand{\UndirectedGraphStruct}[1]{%
  \ensuremath{\mathsf{UGraph}\!\left(#1\right)}%
}
\DeclareRobustCommand{\SimpleUndirectedGraphStruct}[1]{%
  \ensuremath{\mathsf{SGraph}\!\left(#1\right)}%
}
\DeclareRobustCommand{\ConnectedGraphStruct}[1]{%
  \ensuremath{\mathsf{Zus}\!\left(#1\right)}%
}
\DeclareRobustCommand{\AcyclicGraphStruct}[1]{%
  \ensuremath{\mathsf{Kreisfrei}_{\to}\!\left(#1\right)}%
}

% Von einer Kantenrelation abgeleitete Mengen.
\DeclareRobustCommand{\OutNeighbors}[2]{%
  \ensuremath{N^{+}_{#1}\!\left(#2\right)}%
}
\DeclareRobustCommand{\InNeighbors}[2]{%
  \ensuremath{N^{-}_{#1}\!\left(#2\right)}%
}
\DeclareRobustCommand{\GraphNeighbors}[2]{%
  \ensuremath{N_{#1}\!\left(#2\right)}%
}
\DeclareRobustCommand{\ReverseEdges}[1]{%
  \ensuremath{{#1}^{\mathrm{op}}}%
}
\DeclareRobustCommand{\DirectedInducedEdges}[2]{%
  \ensuremath{{#1}[#2]}%
}

% Endliche gerichtete Wege. Der letzte Parameter ist die Kantenlaenge;
% der zugrunde liegende Folgenbereich ist daher NatSeg(n+1).
\DeclareRobustCommand{\DirectedWalk}[4]{%
  \ensuremath{\mathsf{Walk}_{#1,#2}\!\left(#3,#4\right)}%
}
\DeclareRobustCommand{\DirectedPath}[4]{%
  \ensuremath{\mathsf{Path}_{#1,#2}\!\left(#3,#4\right)}%
}
\DeclareRobustCommand{\DirectedCycle}[4]{%
  \ensuremath{\mathsf{Cycle}_{#1,#2}\!\left(#3,#4\right)}%
}
\DeclareRobustCommand{\GraphReach}[1]{%
  \ensuremath{\mathrel{\leadsto_{#1}}}%
}

% ------------------------------------------------------------
% Wälder und Bäume.
% Strukturprädikate erhalten wie die Graphprädikate die gesamte
% Parameterliste als ein Argument. Mengen- und Relationsoperatoren behalten
% getrennte Argumente, etwa \GraphPathsBetween{V}{E}{u}{v}.
% ------------------------------------------------------------
\providecommand{\GraphPathsBetween}[4]{%
  \ensuremath{\mathcal P_{#1,#2}\!\left(#3,#4\right)}%
}
\providecommand{\ForestStruct}[1]{%
  \ensuremath{\mathsf{Wald}\!\left(#1\right)}%
}
\providecommand{\TreeStruct}[1]{%
  \ensuremath{\mathsf{Baum}\!\left(#1\right)}%
}
\providecommand{\PathAcyclicGraphStruct}[1]{%
  \ensuremath{\mathsf{WegAzyk}\!\left(#1\right)}%
}
\providecommand{\RootedTreeStruct}[1]{%
  \ensuremath{\mathsf{WBaum}\!\left(#1\right)}%
}
\providecommand{\TreeParentRel}[5]{%
  \ensuremath{\operatorname{Elter}_{#1,#2,#3}\!\left(#4,#5\right)}%
}
\providecommand{\TreeChildRel}[5]{%
  \ensuremath{\operatorname{Kind}_{#1,#2,#3}\!\left(#4,#5\right)}%
}
\providecommand{\TreeAncestorRel}[5]{%
  \ensuremath{\operatorname{Vorfahr}_{#1,#2,#3}\!\left(#4,#5\right)}%
}
\providecommand{\TreeChildren}[4]{%
  \ensuremath{\operatorname{Kind}_{#1,#2,#3}\!\left[#4\right]}%
}
\providecommand{\TreeLeafPred}[4]{%
  \ensuremath{\operatorname{Blatt}_{#1,#2,#3}\!\left(#4\right)}%
}
\providecommand{\TreeInnerPred}[4]{%
  \ensuremath{\operatorname{Innen}_{#1,#2,#3}\!\left(#4\right)}%
}
\providecommand{\TreeInnerVertices}[3]{%
  \ensuremath{\operatorname{Innen}_{#1,#2,#3}}%
}
\providecommand{\BinaryTreeStruct}[1]{%
  \ensuremath{\mathsf{BinBaum}\!\left(#1\right)}%
}
\providecommand{\FullBinaryTreeStruct}[1]{%
  \ensuremath{\mathsf{VollBinBaum}\!\left(#1\right)}%
}
\providecommand{\OrderedFullBinaryTreeStruct}[1]{%
  \ensuremath{\mathsf{GVollBinBaum}\!\left(#1\right)}%
}

% Hilfspraedikate fuer endliche Wege (auch im Gesamtband verfuegbar).
\providecommand{\FinitePathTransportProperty}[3]{%
  \ensuremath{\mathsf{Transp}_{#1,#2}\!\left(#3\right)}%
}
\providecommand{\FinitePathReversalProperty}[3]{%
  \ensuremath{\mathsf{Umkehr}_{#1,#2}\!\left(#3\right)}%
}
\providecommand{\FinitePathSuffixProperty}[5]{%
  \ensuremath{\mathsf{Schwanz}_{#1,#2,#3,#4}\!\left(#5\right)}%
}
\providecommand{\FinitePathSpliceProperty}[3]{%
  \ensuremath{\mathsf{Füge}_{#1,#2}\!\left(#3\right)}%
}
\providecommand{\FinitePathCycleCore}[3]{%
  \ensuremath{%
    #3\subseteq #1\land #3\neq\varnothing\land
    \forall x\in #3\,\exists y\in #3\,
    \exists m\in\mathbb N\,\exists n\in\mathbb N\,
    \exists p\in\powerset(\mathbb N\times #1)\,
    \exists q\in\powerset(\mathbb N\times #1)\,\Bigl(
      x\neq y\land
      \DirectedPath{#1}{#2}{p}{m}\land p(0)=x\land p(m)=y\land
      \DirectedPath{#1}{#2}{q}{n}\land q(0)=x\land q(n)=y\land
      (m,p)\neq(n,q)\land
      p[\NatSeg{m+1}]\cap q[\NatSeg{n+1}]=\{x,y\}\land
      p[\NatSeg{m+1}]\cup q[\NatSeg{n+1}]\subseteq #3
    \Bigr)%
  }%
}

